% Section 2: Overall Description
\section{Overall Description}

\subsection{Product Perspective}

The Mentor Caring System (MCS) is a new, self-contained software product developed exclusively for the BNBU Mentor Caring Programme, which aims to digitalize and standardize the management of the programme's daily operations and related activities for all undergraduate students of BNBU. This system is an independent supporting management tool and not a follow-on product of any existing software series, nor a replacement for other school information systems.

MCS interfaces with the existing BNBU school database (which stores staff and students' accounts, passwords, names and role attributes) to realize unified user identity authentication. It does not act as a component of a larger system, and all core functions (e.g., mentor allocation, interview record management, appointment arrangement) are built within the independent MCS framework. The system only interacts with the BNBU school database for data reading and verification during user login, without modifying or managing the original data of the school database.

\begin{figure}[htbp]
\centering
\includegraphics[width=0.5\textwidth]{context_model.png}
\caption{Context model of the Mentor Caring System}
\label{fig:context-model}
\end{figure}

\subsection{Product Features}

MCS is built around the core goals of the BNBU Mentor Caring Programme, providing role-based differentiated functional modules to realize the full-process digital management of the programme. The high-level core features are summarized as follows, with detailed specifications in Section 3:

\begin{enumerate}
\item \textbf{Role Management \& Permission Allocation:} Supports hierarchical permission setting for six user roles, realizes unified identity authentication via the BNBU school database, and allows flexible configuration of the school's faculty-department-major organizational architecture (manual/Excel import).

\item \textbf{Mentor \& Student Allocation Management:} Facilitates batch import of mentor-student allocation information via Excel, supports manual adjustment of group mentors and group student members, and retains all historical information of students during any adjustment operation.

\item \textbf{Information Query \& Statistics:} Provides multi-dimensional search functions for student, mentor and system log information (supporting precise and conditional screening), and enables batch export of specified information to Word files based on academic year, department, mentor and student name.

\item \textbf{Interview Management:} Supports mentors to add, modify and delete interview records of their supervised students, and allows batch export of interview records to Word files for statistical and archival purposes.

\item \textbf{Online Communication \& Appointment:} Realizes two-way text communication between different roles, and supports the mentor-student interview appointment function (including mentor-set available time slots, venue confirmation and email notification for both parties).

\item \textbf{Special Case Processing:} Establishes a tiered forwarding mechanism for student special cases, from mentors to MCP coordinators and then to faculty consultants, to ensure timely follow-up of key student issues.

\item \textbf{System Log Monitoring:} Records the login/logout time and core function usage of all users, and provides log query and viewing permissions for designated roles to track system operation status.
\end{enumerate}

The Mentor Caring System (MCS) provides a centralized digital platform to support the BNBU Mentor Caring Programme. The system organizes its major functions into several groups of related capabilities.

First, the system supports role-based access control and identity authentication through the BNBU school database, ensuring that administrators, faculty consultants, mentors, students, MCP coordinators and supporting staff can access functions according to their permissions.

Second, the system manages mentor--student allocation and organizational structures, allowing faculty consultants to import or modify information about faculties, departments, majors, mentoring groups and student membership through Excel templates or manual operations.

Third, the system provides information management and query services, enabling users to search student records, mentor information and system logs, and export selected records into Word documents for reporting and archiving.

Fourth, the system supports interview record management and appointment scheduling between mentors and students. Mentors can create, modify or delete interview records and arrange available appointment time slots for students.

Finally, the system enables internal communication among different user roles through text messages and notifications, including forwarding special student cases from mentors to MCP coordinators and faculty consultants for further follow-up.

\begin{figure}[htbp]
\centering
\includegraphics[width=\textwidth]{usecase_fig1.png}
\caption{Use Case Diagram: Mentor, Student and Coordinator}
\label{fig:usecase-fig1}
\end{figure}

\begin{figure}[htbp]
\centering
\includegraphics[width=\textwidth]{usecase_fig2_fig3.png}
\caption{Use Case Diagram: Administrator, Faculty Consultant and Supporting Staff}
\label{fig:usecase-fig2-fig3}
\end{figure}

\subsection{User Classes and Characteristics}

MCS is designed for six distinct user classes, differentiated by security/privilege levels, subset of product functions used and frequency of use, all of whom access the system via BNBU's unified account authentication. Core user classes and their pertinent characteristics are as follows, with faculty consultants and mentors as the favored user classes (the main daily operators of the system):

\begin{enumerate}
\item \textbf{Administrator}
\begin{itemize}
\item Privilege Level: Highest system privilege; only 1 user for the entire university.
\item Core Characteristics: Non-technical or technical background (BNBU school management staff); low frequency of system use (mainly for initial setup and occasional maintenance).
\item Key Functions: Designate/modify/delete faculty consultants, configure organizational architecture (faculty-department-major), set up supporting staff.
\item Unique Constraints: Cannot act as a mentor in the system.
\end{itemize}

\item \textbf{Faculty Consultant}
\begin{itemize}
\item Privilege Level: High; one or more per BNBU faculty (FST, FHSS, FSM, DCC).
\item Core Characteristics: Faculty management staff with familiarity with Excel operation; medium-to-high frequency of use (key operation at the start of each academic year, regular daily maintenance).
\item Key Functions: The core manager of the faculty-level MCS, responsible for mentor-student allocation import, group adjustment, MCP coordinator designation, multi-dimensional information search/export and cross-role communication.
\item Favored User Class: Primary daily operator of the system, core to the normal operation of the Mentor Caring Programme.
\end{itemize}

\item \textbf{Mentor}
\begin{itemize}
\item Privilege Level: Medium; assigned to mentor groups for undergraduate students.
\item Core Characteristics: BNBU teaching staff; high frequency of use (daily communication with students, interview record management, regular appointment arrangement).
\item Key Functions: Restricted to operating on their own supervised groups, including student information query, interview record editing, mentor-student communication/appointment, special case forwarding and interview record export.
\item Favored User Class: The most frequent operator of the system, the direct executor of the Mentor Caring Programme's core work.
\end{itemize}

\item \textbf{Student}
\begin{itemize}
\item Privilege Level: Low; all BNBU undergraduate students in mentor groups.
\item Core Characteristics: Non-technical background; variable use frequency (mainly for checking mentor information, responding to messages/appointments).
\item Key Functions: View personal mentor's basic information, two-way communication and interview appointment with the mentor.
\item Unique Constraints: Cannot access the system if not assigned to any mentor group (system displays a warning and forces logout).
\end{itemize}

\item \textbf{MCP Coordinator}
\begin{itemize}
\item Privilege Level: Medium; one per BNBU department, designated by faculty consultants.
\item Core Characteristics: Department management staff; medium frequency of use (mainly for special case forwarding and cross-role communication).
\item Key Functions: Search mentor/student information, forward student special cases from mentors to faculty consultants, and two-way communication with students, faculty consultants and mentors.
\item Flexible Attribute: May concurrently hold the role of a mentor for a mentor group.
\end{itemize}

\item \textbf{Supporting Staff}
\begin{itemize}
\item Privilege Level: Low; designated by the administrator (optional role).
\item Core Characteristics: School administrative staff; low frequency of use (auxiliary system management).
\item Key Functions: View all students' log information, respond to user feedback on system usage (the latter function is non-mandatory).
\end{itemize}
\end{enumerate}

\subsection{Operating Environment}

\textbf{Server-side Environment}

The Mentor Caring System (MCS) shall be deployed on a standard enterprise-level server capable of hosting a web application and a dedicated relational database. The server must maintain a stable network connection to the BNBU school database server for user identity verification and to the BNBU email server for automatic notification delivery. The specific server hardware configuration and operating system shall be determined during the system design phase based on the expected scale of university-wide usage.

\textbf{Recommended Configuration}

Hardware Platform: Enterprise-level server (CPU: Intel Xeon Silver/Gold or equivalent or above; RAM: 32GB or above; Hard Disk: 1TB SSD or above) with high-speed and stable network connectivity.

Operating System: Enterprise-grade Linux distribution (such as Rocky Linux 8/9 or Ubuntu Server 20.04/22.04 LTS) or Windows Server 2022 (64-bit).

Supporting Software: Relational database management system (MySQL 8.0 latest stable version / SQL Server 2019 or above); Web server (Nginx 1.20 or above / Apache Tomcat 9.0 or above).

\textbf{Client-side Environment}

Users shall access the system through standard personal computers (desktops or laptops) equipped with mainstream modern web browsers. Client devices must have basic network access to the BNBU internal campus network, or external network access through the university's VPN service. Users are also required to have Microsoft Excel (for importing and editing system-related data files as specified in Appendix A--D) and Microsoft Word (for viewing and editing exported report files) installed on their devices.

\textbf{Recommended Configuration}

Hardware Platform: Personal computer with Intel i5 / AMD Ryzen 5 or above, RAM 8GB or above, SSD storage and stable broadband network connection.

Operating System: Windows 11 (64-bit) or macOS 13 or above.

Browser: Latest stable versions of Chrome, Edge or Firefox.

Supporting Software: Microsoft Office 2019 / Microsoft 365 or above.

\textbf{Coexistence Requirements}

The system must peacefully coexist with the existing BNBU school database system and the school's email server (for sending system message notifications), without occupying excessive system resources or affecting the normal operation of other BNBU internal information systems.

\subsection{Design and Implementation Constraints}

The design and implementation of MCS are subject to the following constraints, which limit the development options for the development team:

\begin{enumerate}
\item \textbf{Functional Adaptability Constraint:} The system must support the dynamic addition of new faculties in BNBU without modifying the core code, to adapt to the school's future organizational structure adjustments.
\item \textbf{Data Retention Constraint:} All student information (including personal data, interview records, group history) must be permanently retained in the system; no data deletion is allowed for any operation (e.g., removing a student from a group, changing a group's mentor).
\item \textbf{Data Interaction Constraint:} The system can only read data from the BNBU school database (for user authentication) and cannot modify, delete or store any original data of the school database; all system-specific data (e.g., mentor-student allocation, interview records) must be stored in MCS's independent database.
\item \textbf{File Format Constraint:} The system must strictly support the Excel/Word file formats specified in the BNBU Mentor Caring Programme appendices (Appendix A-D) for data import and export; no custom non-standard file formats are allowed.
\item \textbf{Message Notification Constraint:} All system internal messages (e.g., communication, appointment) must trigger an automatic email notification to the receiver via the BNBU email server; receivers can only view and respond to messages after logging into the system.
\item \textbf{Appointment Time Constraint:} The mentor-student interview appointment function must set available time slots in 30-minute fixed intervals; no custom time slot durations are supported.
\item \textbf{Permission Isolation Constraint:} Mentors can only access and operate data related to their own supervised groups; no cross-group data access is allowed, and the system must strictly enforce role-based access control (RBAC) for all user classes.
\item \textbf{Maintainability Constraint:} The system code must follow standard software development programming norms (e.g., modular design, detailed code comments) to facilitate subsequent maintenance and version iteration by BNBU's technical staff.
\item \textbf{Excel Import Constraint:} When importing Excel files into the system, existing information in the system must not be overwritten; only non-existing information is updated, and the import logic must follow this rule for all Excel-based data operations.
\end{enumerate}

\subsection{User Documentation}

The following user documentation components will be delivered along with the MCS software, all complying with BNBU's internal information system documentation standards and the general software user manual writing norms. All documents will be provided in both electronic (PDF format) and editable (Word format) versions, and published on the BNBU internal system platform for easy access by all users:

\begin{enumerate}
\item \textbf{MCS Quick Start Guide:} A concise document for all user classes, including basic system access steps, account authentication instructions, core function operation shortcuts and common problem troubleshooting (1-2 pages for quick reference).

\item \textbf{MCS Detailed User Manual:} A comprehensive document classified by user roles, including detailed step-by-step operation instructions for all functions, screenshot demonstrations of key operations, file format specifications (Appendix A-D) and system error code explanations. Separate chapters are provided for each role to ensure targeted readability.

\item \textbf{MCS Administrator Maintenance Manual:} A dedicated document for the system administrator, including organizational architecture configuration, role permission setting, supporting staff management, system log viewing and basic system maintenance operations (e.g., data backup, Excel template update).

\item \textbf{MCS File Template Package:} A compressed package containing the standard Excel templates for Appendix A (organizational architecture), Appendix B (mentor-student allocation) and Appendix C (MCP coordinator designation), with built-in data validation rules to ensure user input compliance.
\end{enumerate}

\subsection{Assumptions and Dependencies}

\textbf{Assumptions}

The following assumed factors (not yet verified as absolute facts) affect the requirements and implementation of MCS; if these assumptions are incorrect, unshared or changed, the project development and system operation will be impacted:

\begin{enumerate}
\item The BNBU school database is permanently available and stably connected during the entire operation period of MCS, and provides a fixed data interface for MCS to read user account, password, name and role information.
\item The BNBU email server is always operational and supports the system's automatic email notification function; the server has no sending frequency or content restrictions for internal school email notifications.
\item All MCS users (especially faculty consultants and mentors) have basic computer operation skills and are proficient in using Microsoft Excel and Word (the core office software for the system).
\item BNBU will not adjust the core attributes of the user account (e.g., student ID, staff email) during the system operation period; any account attribute adjustment will be notified to the development team in advance.
\item The non-mandatory function of ``supporting staff responding to user feedback'' will not require additional integration with BNBU's feedback system; the supporting staff only needs to process feedback via manual record within MCS.
\end{enumerate}

\textbf{Dependencies}

The development and operation of MCS have the following key external dependencies, which are essential for the realization of core system functions:

\begin{enumerate}
\item \textbf{BNBU School Database Interface:} The system is dependent on the provision of a valid, read-only data interface by the BNBU Information Technology Center; without this interface, the unified user identity authentication function cannot be realized.
\item \textbf{BNBU Email Server Service:} The system is dependent on the BNBU email server to provide SMTP/POP3 services for automatic message notification; the interruption of this service will result in the failure of the email notification function for communication and appointments.
\item \textbf{Microsoft Office Software:} The system's data import/export function is dependent on the mainstream Microsoft Office (Excel/Word) software used by all BNBU users; the system does not support other office software with incompatible file formats.
\item \textbf{BNBU Internal Network Environment:} The system's core data interaction (e.g., accessing the school database) is dependent on the stability of the BNBU internal campus network; external network access requires the school's VPN service to ensure secure and stable connection to the system server.
\end{enumerate}
